\chapter*{Abstract}
\addcontentsline{toc}{chapter}{Abstract}

Ordered guard specifications say which scenario must fire first when several
guards hold.
Large language models rarely obey that rule: they look right on everyday
inputs and fail on the thin overlaps that formal contracts are meant to pin
down.
Unit tests almost never sample those regions, so test-based repair leaves the
defects in place.
That failure class is \emph{specification-conformance defects}.

The key insight is to treat the specification itself as the evidence source.
The \textbf{Specification-guided Defect Prevention Framework} (\textsc{SgDP}),
instantiated as \textbf{HSP-Agile} for SOFL/FSF, builds first-match SMT
witnesses, turns each failure into a typed repair record, and releases a
candidate only when every witness matches and no high-severity pattern remains.

On a 120-task precedence-sensitive benchmark, HSP-Agile reaches 90.4\% mean
strict formal conformance (+6.2\,pp over one-shot LLM, +2.4\,pp over
test-feedback).
The advantage grows with guard overlap; SMT keeps boundary coverage where
random sampling fails; the typed repair record alone accounts for +7.7\,pp
over test-only prompts (E6).
The same loop transfers to Mini-Z and Mini-StateMachine with adapter-level
changes.
Tasks, logs, and scripts are released.

\paragraph{Keywords.}
Specification-guided defect prevention; formal specification; LLM synthesis; SOFL;
counterexample-guided repair; SMT verification; mutation testing.
